\documentclass[11pt,a4paper]{article}
% ---------------------------------------------------------------------------
%  FermiWick -- common style of the showcase and of the manual.
%  The examples are computed by Development/run_examples.wl: examples/NAME.wl is the code, NAME.tex the output as a
%  formula, NAME.txt the output as text.
% ---------------------------------------------------------------------------
\usepackage[utf8]{inputenc}
\usepackage[T1]{fontenc}
\usepackage[sc]{mathpazo}
\usepackage[scaled=0.88]{helvet}
\renewcommand{\ttdefault}{txtt}
\usepackage{amsmath,amssymb,bm}
\allowdisplaybreaks[2]
\usepackage[a4paper,margin=2.2cm,top=2.4cm,bottom=2.4cm,headheight=22pt]{geometry}
\usepackage{microtype}
\usepackage[dvipsnames]{xcolor}
\usepackage{graphicx}
\usepackage{listings}
\usepackage{simpler-wick}
\usepackage[most]{tcolorbox}
\usepackage{booktabs,array,tabularx}
\usepackage{titlesec}
\usepackage{enumitem}
\usepackage{fancyhdr}
\usepackage{needspace}
\usepackage{etoolbox}
\usepackage[colorlinks=true,linkcolor=accent,urlcolor=accent,citecolor=accent]{hyperref}

\definecolor{accent}{HTML}{8E2A7B}
\definecolor{accentlight}{HTML}{F1DCEC}
\definecolor{accentpale}{HTML}{FAF0F7}
\definecolor{codebg}{HTML}{FBF6FA}
\definecolor{codeframe}{HTML}{E1CCDC}
\definecolor{keyword}{HTML}{8E2A7B}
\definecolor{comment}{HTML}{7D6A78}
\definecolor{checkgreen}{HTML}{2E7D32}

\titleformat{\section}{\sffamily\Large\bfseries\color{accent}}{\thesection}{0.8em}{}
\titleformat{\subsection}{\sffamily\large\bfseries\color{accent}}{\thesubsection}{0.7em}{}
\titleformat{\paragraph}[runin]{\sffamily\bfseries\color{accent}}{}{0pt}{}[.]
\titlespacing*{\section}{0pt}{2.6ex plus 1ex minus 0.4ex}{1.2ex}
\titlespacing*{\subsection}{0pt}{2.2ex plus 0.6ex minus 0.3ex}{0.8ex}
\titlespacing*{\paragraph}{0pt}{1.2ex plus 0.4ex}{0.6em}
\makeatletter
\pretocmd{\section}{\needspace{8\baselineskip}}{}{}
\pretocmd{\subsection}{\if@nobreak\else\needspace{6\baselineskip}\fi}{}{}
\makeatother
\setcounter{tocdepth}{2}
\setlist{itemsep=1pt, topsep=3pt}
\setlength{\emergencystretch}{1.5em}
\setlength{\parindent}{0pt}
\setlength{\parskip}{4pt plus 1.5pt minus 1pt}

\pagestyle{fancy}
\fancyhf{}
\fancyhead[R]{\sffamily\small\color{comment}\nouppercase{\rightmark}}
\fancyfoot[C]{\sffamily\small\thepage}
\renewcommand{\sectionmark}[1]{\markboth{\thesection\enspace #1}{\thesection\enspace #1}}
\renewcommand{\headrulewidth}{0.4pt}
\renewcommand{\headrule}{\hbox to\headwidth{\color{codeframe}\leaders\hrule height \headrulewidth\hfill}}

\lstdefinelanguage{WL}{
  morekeywords={NormalOrder,OperatorCommutator,OperatorAnticommutator,ToSpinChargeBasis,ToLadderBasis,GroupLocalOperators,
    OperatorCoefficients,CleanDeltas,ResetNormalOrder,DefineOperator,QuasiFreeState,SpinSingletState,GeneralState,StateSelect,
    LocalPairQ,AnomalousPairQ,Contraction,WickVEV,VerifyNormalOrder,MeanField,DefinePeierlsHopping,DefineFlux,
    GutzwillerProjector,LocalSandwich,GutzwillerVEV,GutzwillerSimplify,SdotS,SdotScrossS,SpinOperator,EmptyProjector,
    SingleOccupationProjector,cDag,nOp,nTotOp,SzOp,SpOp,SmOp,DOp,Bdag,
    FockSpace,FockSpaceQ,FockSector,FockSectorQ,ModeIndex,OccupationBasis,BasisState,StateToKet,FockMatrix,FockExpectation,
    SpinSquared,ParticleHoleOperator,TimeReversal,SectorBlock,LowestStates,LehmannGreen,GreenMatrix,ClearFockCache},
  sensitive=true, morecomment=[n]{(*}{*)}, morestring=[b]",
}
\lstset{language=WL, basicstyle=\ttfamily\small, keywordstyle=\color{keyword}\bfseries, commentstyle=\color{comment}\itshape,
  stringstyle=\color{BrickRed}, columns=fullflexible, keepspaces=true, breaklines=true, breakatwhitespace=true,
  breakindent=2em, showstringspaces=false, upquote=true, extendedchars=true,
  literate={·}{{$\cdot$}}1 {φ}{{$\varphi$}}1 {Φ}{{$\Phi$}}1 {σ}{{$\sigma$}}1 {ρ}{{$\rho$}}1 {κ}{{$\kappa$}}1 {↑}{{$\uparrow$}}1 {↓}{{$\downarrow$}}1 {χ}{{$\chi$}}1 {Δ}{{$\Delta$}}1 {θ}{{$\theta$}}1 {ν}{{$\nu$}}1 {η}{{$\eta$}}1}
\lstdefinestyle{plainoutput}{language={}, basicstyle=\ttfamily\small, breaklines=true, columns=fullflexible, keepspaces=true}

\newcommand{\pkg}[1]{\texttt{#1}}
\newcommand{\fn}[1]{\texttt{\color{accent}#1}}
\newcommand{\ok}{\textcolor{checkgreen}{\ensuremath{\checkmark}}}
\newcommand{\exmissing}[1]{\textcolor{red}{[missing example: \detokenize{#1}]}}
\makeatletter
\newcommand{\excodeafter}{\tcb@parfillskip@check\par\ifvmode\nobreak\@nobreaktrue\nointerlineskip\addvspace{\glueexpr2pt-\parskip}\fi}
\makeatother
\newcommand{\excode}[1]{\ifhmode{\widowpenalty=10000\par}\nobreak\fi\IfFileExists{examples/#1.wl}{\tcbinputlisting{listing only,
  colback=codebg, colframe=codeframe, boxrule=0.5pt, arc=2pt, left=6pt, right=6pt, top=3pt, bottom=3pt, before skip=4pt,
  after=\excodeafter, listing file={examples/#1.wl}, listing options={language=WL}}}{\exmissing{#1}}}
\newtcolorbox{exoutput}{colback=white, colframe=codeframe, boxrule=0.5pt, arc=2pt, left=24pt, right=6pt, top=3pt, bottom=3pt,
  before skip=2pt, after skip=8pt, breakable, overlay unbroken and first={\node[anchor=north west,
  font=\scriptsize\sffamily\color{comment}] at ([xshift=4pt, yshift=-5pt]frame.north west) {out};}}
\newcommand{\exformula}[1]{\begin{exoutput}\raggedright\IfFileExists{examples/#1.tex}{$\displaystyle\input{examples/#1.tex}$}{\exmissing{#1}}\end{exoutput}}
\newcommand{\extext}[1]{\IfFileExists{examples/#1.txt}{\tcbinputlisting{listing only, colback=white, colframe=codeframe,
  boxrule=0.5pt, arc=2pt, left=24pt, right=6pt, top=2pt, bottom=2pt, before skip=2pt, after skip=8pt,
  listing file={examples/#1.txt}, listing options={style=plainoutput}, overlay unbroken and first={\node[anchor=north west,
  font=\scriptsize\sffamily\color{comment}] at ([xshift=4pt, yshift=-5pt]frame.north west) {out};}}}{\exmissing{#1}}}
\newcommand{\example}[1]{\excode{#1}\IfFileExists{examples/#1.tex}{\exformula{#1}}{\IfFileExists{examples/#1.txt}{\extext{#1}}{}}}

\newcommand{\titleblock}[2]{%
  \begin{center}
  \includegraphics[width=2.2cm]{../Logo/FermiWick_logo.pdf}\\[6pt]
  {\sffamily\Huge\bfseries\color{accent}FermiWick}\\[4pt]
  {\sffamily\Large #1}\\[8pt]
  {\small #2}\\[2pt]
  {\small Nazzareno Africani (SISSA) \textperiodcentered{} version 0.2.0 \textperiodcentered{}
   \href{https://github.com/n4sd4q/FermiWick}{github.com/n4sd4q/FermiWick}}
  \end{center}\vspace{6pt}}
\newtcblisting{code}{listing only, colback=codebg, colframe=codeframe, boxrule=0.5pt, arc=2pt, left=6pt, right=6pt,
  top=3pt, bottom=3pt, listing options={language=WL}}
\newcommand{\extypeset}[1]{\begin{exoutput}\centering\small\IfFileExists{examples/#1_typeset.tex}{$\displaystyle\input{examples/#1_typeset.tex}$}{\exmissing{#1}}\end{exoutput}}

\hypersetup{pdftitle={FermiWick: reference manual},pdfauthor={Nazzareno Africani}}
\fancyhead[L]{\sffamily\small\color{accent}FermiWick \textperiodcentered{} reference manual}
\newcommand{\fnentry}[2]{\par\needspace{4\baselineskip}\noindent\fn{#1}\par\nobreak
  \noindent\begin{minipage}[t]{\linewidth}\leftskip=1.2em #2\end{minipage}\par}

\begin{document}
\titleblock{Reference manual}{Every function, option and convention, with an example computed by the package}

{\small\tableofcontents}

\section{Installation and loading}
FermiWick is a paclet. Install it once from its folder (the one with \texttt{PacletInfo.wl}), then load it in every
session, with the Gutzwiller and Fock modules when needed:
\begin{code}
PacletInstall[CreatePacletArchive["/path/to/FermiWick/FermiWick", $TemporaryDirectory]]
Needs["FermiWick`"]; Needs["FermiWick`Gutzwiller`"]; Needs["FermiWick`Fock`"];
\end{code}
To use it without installing, \texttt{PacletDirectoryLoad["/path/to/FermiWick/FermiWick"]} before \texttt{Needs}. The
guided tour, \texttt{Examples.nb} in the repository, runs the examples of the showcase cell by cell.
FermiWick has no dependencies. It is tested with Mathematica 15.0. The test suite runs with
\texttt{math -script Tests/RunTests.wls}.

\section{Conventions}
\begin{itemize}
\item \textbf{Operators.} \fn{c[l1, l2, \dots]} and \fn{cDag[l1, l2, \dots]} annihilate and create a fermion in the mode
  with labels \texttt{l1, l2, \dots}; \fn{DeclareModeLabels} sets the meaning of each label. By default \fn{c[i]} is a
  spinless fermion on site \texttt{i} and \fn{c[i, s]} a fermion on site \texttt{i} with spin \texttt{s}. A site is any
  expression, and two different symbolic sites are different sites: \texttt{KroneckerDelta[i, j]} is $0$. A spin is
  \texttt{SpinUp}, \texttt{SpinDown} or a symbol, split into $\uparrow$ and $\downarrow$ when needed. Generic labels (orbitals,
  momenta, \dots) are compared with \texttt{KroneckerDelta}, which is kept and substituted in what it multiplies.
\item \textbf{Products.} \texttt{a ** b} is the product in the given order; \texttt{a \textperiodcentered{} b} (\texttt{\textbackslash[CenterDot]})
  is the product brought to normal order. Scalars (anything free of operators) factor out.
\item \textbf{Normal order.} All $c^\dagger$ to the left of all $c$; results are expanded sums of monomials.
\item \textbf{Contractions.} On a pair of sites, in spin channels:
  $\langle c^\dagger_{i\sigma}c_{j\sigma'}\rangle=\chi_{ij}\delta_{\sigma\sigma'}+\sum_a\chi^a_{ij}\sigma^a_{\sigma'\sigma}$ and
  $\langle c_{i\sigma}c_{j\sigma'}\rangle=[(\Delta_{ij}+\sum_a\Delta^a_{ij}\sigma^a)\epsilon]_{\sigma\sigma'}$, with $\epsilon_{\uparrow\downarrow}=1$,
  $a=x,y,z$ (written \texttt{"x"}, \texttt{"y"}, \texttt{"z"}): \fn{\textbackslash[Chi][i, j]}, \fn{\textbackslash[Chi][i, j, a]},
  \fn{\textbackslash[CapitalDelta][i, j]}, \fn{\textbackslash[CapitalDelta][i, j, a]}. $\Delta$ is symmetric and $\Delta^a$
  antisymmetric in the sites; Hermiticity gives $\chi_{ji}=\chi^*_{ij}$, $\chi^a_{ji}=\chi^{a*}_{ij}$, and the amplitudes are
  written with the sites in canonical order (\texttt{Conjugate[\textbackslash[Chi][1, 2]]} for $\chi_{21}$). On one site $\chi_{ii}$
  is the density per spin and $\chi^a_{ii}=\langle S^a_i\rangle$. In the general state by spin components,
  $\rho_{i\sigma j\sigma'}=\langle c^\dagger_{i\sigma}c_{j\sigma'}\rangle$ and $\kappa_{i\sigma j\sigma'}=\langle c_{i\sigma}c_{j\sigma'}\rangle$.
\item \textbf{Complex conjugation.} The amplitudes are complex; every other symbol is taken as real by \fn{Dagger}, unless
  declared complex.
\end{itemize}
Symbolic spins and sites:
\example{symbolicspin}
\example{samesite}

\section{Operators}
\fnentry{c, cDag}{The ladder operators. Their square vanishes:}
\example{pauli}
\fnentry{DeclareModeLabels, \$ModeLabels}{\fn{DeclareModeLabels[\{type1, type2, \dots\}]} declares the labels of the
operators with that number of labels: \texttt{"Site"} (locality: products are reduced site by site), \texttt{"Spin"} (a spin
$\frac12$; the spin operators below need modes \fn{c[i, s]}), or any other name, a generic label. Defaults:
\texttt{\{"Site"\}}, \texttt{\{"Site", "Spin"\}}, and \texttt{\{"Site", "Label", \dots\}} for more labels. \fn{\$ModeLabels} holds
the declarations. Modes with no site are normal-ordered without the site-by-site reduction.}
\example{labels}
\fnentry{SpinUp, SpinDown, \$Spins}{\fn{SpinUp} and \fn{SpinDown} are the spins $\uparrow$ and $\downarrow$, \fn{\$Spins} their
list, for spin sums:}
\example{spins}
\fnentry{nOp, nTotOp, SzOp, SpOp, SmOp, DOp, B, Bdag}{The local operators of a site with spin; they stay compact, and
\fn{ToLadderBasis} writes them with $c$, $c^\dagger$. All the composite operators and their definitions:}
% The composite operators of FermiWick and their definitions (included by the showcase and by the manual).
\begin{center}\small
\renewcommand{\arraystretch}{1.25}
\begin{tabular}{lll}\toprule
FermiWick & operator & definition \\ \midrule
\texttt{nOp[i, s]} & $n_{i\sigma}$ & $c^\dagger_{i\sigma}c_{i\sigma}$ \\
\texttt{nTotOp[i]}, \texttt{nOp[i]} & $n_i$ & $n_{i\uparrow}+n_{i\downarrow}$ \\
\texttt{SzOp[i]} & $S^z_i$ & $\tfrac12\,(n_{i\uparrow}-n_{i\downarrow})$ \\
\texttt{SpOp[i]}, \texttt{SmOp[i]} & $S^+_i$, $S^-_i$ & $c^\dagger_{i\uparrow}c_{i\downarrow}$, \ $c^\dagger_{i\downarrow}c_{i\uparrow}$ \\
\texttt{SpinOperator[i, a]} & $S^a_i$ & $\tfrac12\sum_{\sigma\sigma'}c^\dagger_{i\sigma}\,\sigma^a_{\sigma\sigma'}\,c_{i\sigma'}$, \ $a=1,2,3$ \\
\texttt{DOp[i]} & $D_i$ (double occupancy) & $n_{i\uparrow}n_{i\downarrow}$ \\
\texttt{Bdag[i]}, \texttt{B[i]} & $B^\dagger_i$, $B_i$ (doublon) & $c^\dagger_{i\uparrow}c^\dagger_{i\downarrow}$, \ $c_{i\downarrow}c_{i\uparrow}$ \\
\texttt{EmptyProjector[i]} & $P^0_i$ & $(1-n_{i\uparrow})(1-n_{i\downarrow})=1-n_i+D_i$ \\
\texttt{SingleOccupationProjector[i]} & $P^1_i$ & $n_i-2D_i$ \\
\texttt{GutzwillerProjector[i]} & $P^G_i$ & $\sqrt2\,[\cos\theta\,P^1_i+\sin\theta\,(P^0_i+D_i)]$ \\
\texttt{SdotS[i, j]} & $\mathbf S_i\cdot\mathbf S_j$ & $S^z_iS^z_j+\tfrac12(S^+_iS^-_j+S^-_iS^+_j)$ \\
\texttt{SdotScrossS[i, j, k]} & $\mathbf S_i\cdot(\mathbf S_j\times\mathbf S_k)$ & $\sum_{abc}\epsilon_{abc}\,S^a_iS^b_jS^c_k$ \\ \bottomrule
\end{tabular}
\end{center}
Products of these on one site reduce with the algebra of the fermions (e.g.\ $D_in_i=2D_i$, $(S^z_i)^2=\tfrac14n_i-\tfrac12D_i$);
operators on several sites and products with symbolic sites stay as products. \texttt{GutzwillerProjector} is in
\texttt{FermiWick`Gutzwiller`}; one-site operators on several sites are products, \texttt{DOp[i, j]} $=D_iD_j$.

\example{composite}
\fnentry{SpinOperator, SdotS, SdotScrossS}{\fn{SpinOperator[i, a]} is $S^a_i$, $a=1,2,3$; \fn{SdotS[i, j]} is
$\mathbf S_i\cdot\mathbf S_j$ and \fn{SdotScrossS[i, j, k]} the scalar chirality $\mathbf S_i\cdot(\mathbf S_j\times\mathbf S_k)$.}
\example{spincomponents}
\fnentry{EmptyProjector, SingleOccupationProjector}{$P^0_i=1-n_i+D_i$ and $P^1_i=n_i-2D_i$, so that $P^0+P^1+D=1$; with
several sites they are products.}
\example{projectors}
\fnentry{DefineOperator}{\fn{DefineOperator[head, n, definition, display]} defines an operator on $n$ sites: \texttt{definition}
is a function of the sites written with known operators, \texttt{display} a function of the sites that gives its form on
screen. The operator stays compact; \fn{NormalOrder}, \fn{WickVEV} and \fn{MeanField} expand it, and \fn{ToSpinChargeBasis}
recognizes it again. Options: \texttt{"Recognize"} (default \texttt{True} for $n\ge2$) and \texttt{"Required" -> All}
(recognize it only when it is present whole; by default its diagonal part may leave a remainder, as in
$J\,\mathbf S\cdot\mathbf S+(J_z-J)S^zS^z$). A one-site operator on several sites is the product. \fn{\$DefinedOperators} is the
registry.}
\example{define}

\section{Products and bases}
\fnentry{NormalOrder}{Brings an expression to normal order. Products are separated site by site: each factor acts on one
site and is even or odd, so the sites are reordered with the fermion sign only and each site is reduced once.}
\fnentry{CenterDot}{\texttt{a \textperiodcentered{} b \textperiodcentered{} \dots} is the normal-ordered product.}
\example{centerdot}
\fnentry{OperatorCommutator, OperatorAnticommutator}{$ab-ba$ and $ab+ba$, in normal order.}
\example{anticommutator}
\fnentry{ToSpinChargeBasis}{Writes an expression with $n$, $S^z$, $S^\pm$, $D$, $B$, $B^\dagger$ and recognizes the defined
operators; \texttt{"Recognize" -> False} leaves them expanded.}
\example{tospincharge}
\example{norecognize}
\fnentry{ToLadderBasis}{Writes the composite and the defined operators with $c$ and $c^\dagger$.}
\fnentry{GroupLocalOperators}{Normal-orders and groups the ladder operators site by site, with the fermion sign.}
\example{group}
\fnentry{OperatorCoefficients}{The association monomial $\to$ coefficient of an expression in normal order; the constant
has the key \texttt{1}. \fn{OperatorCoefficients[expr, m]} is the coefficient of one monomial, written in any order. It reads
the fields of a mean-field Hamiltonian.}
\example{coefficients}
\example{coefficient}
\fnentry{Dagger}{The Hermitian conjugate: products reversed, $c\leftrightarrow c^\dagger$, $S^+\leftrightarrow S^-$,
$B\leftrightarrow B^\dagger$, coefficients conjugated. The amplitudes $\chi$, $\Delta$, $\rho$, $\kappa$ are complex, the other
symbols real; option \texttt{"Complex" -> \{\dots\}} declares more complex symbols. Operators defined with
\fn{DefineOperator} are recognized as Hermitian when they are; otherwise they are expanded.}
\example{dagger}
\fnentry{PlusHC, HermitianOperatorQ}{\fn{PlusHC[expr]} is $\texttt{expr}+\texttt{expr}^\dagger$; \fn{HermitianOperatorQ[expr]}
checks $\texttt{expr}=\texttt{expr}^\dagger$. Options as \fn{Dagger}.}
\example{plushc}
\fnentry{CleanDeltas, ResetNormalOrder, \$SpinSymbols}{\fn{CleanDeltas} simplifies the \texttt{KroneckerDelta} of sites and
spins and does the sums over the spin; \fn{ResetNormalOrder[]} clears the caches; \fn{\$SpinSymbols} lists the symbols
used as spins.}

\section{Gaussian states}
A Gaussian state is an object that gives the contraction $\langle O_1O_2\rangle$ of two ladder operators. It is declared by
its symmetries, or by rules for the contractions. \fn{state["Symmetries"]}, \fn{state["Description"]} and
\fn{state["Contraction"]} give its parts.
\fnentry{GaussianState}{Gaussian states, in which every expectation value follows from the contractions of two operators,
are also called quasi-free states (pure, as Slater or BCS states, or mixed, as Gibbs states of quadratic Hamiltonians).
\fn{GaussianState[symmetries, options]}: the state with the given symmetries, written in spin
channels (Section~2). \fn{GaussianState[]} is the general state. Each symmetry sets channels to zero or makes them real:
\begin{tabular}{ll}\toprule symmetry & what it leaves \\ \midrule
\texttt{"SpinRotation"} & $\chi$, $\Delta$ \\
\texttt{"SpinRotationZ"} & $\chi$, $\chi^z$, $\Delta$, $\Delta^z$ \\
\texttt{"Charge"} & $\chi$, $\chi^a$ (no pairing) \\
\texttt{"TimeReversal"} & $\chi$, $\Delta$ real; $\chi^a$, $\Delta^a$ imaginary; no magnetization \\ \bottomrule\end{tabular}

Options:
\begin{tabular}{lll}\toprule option & default & \\ \midrule
\texttt{"Hopping"}, \texttt{"Pairing"} & $\chi$, $\Delta$ & names of the amplitudes \\
\texttt{"Filling"} & \texttt{Automatic} & $\chi_{ii}$; \texttt{Automatic}: the symbol $\chi_{ii}$ \\
\texttt{"Magnetization"} & \texttt{Automatic} & $\langle S^z_i\rangle$, a vector, or a function of the site \\
\texttt{"OnSitePairing"} & \texttt{Automatic} & $\Delta_{ii}$ \\
\texttt{"Hermitian"} & \texttt{True} & $\chi_{ji}=\chi^*_{ij}$; \texttt{False}: $\chi_{ji}$ is a symbol of its own \\
\texttt{"Rules"} & \texttt{\{\}} & replacements applied to the results \\ \bottomrule\end{tabular}

\fn{GaussianState[\{\{op1, op2\} :> value, \dots\}]} declares the contractions by rules instead, with patterns in the
operators. Pairs with no rule contract to zero, and $\langle c\,c^\dagger\rangle$ follows from $\langle c^\dagger c\rangle$ by the
anticommutator when it has no rule of its own. Two sites with hopping $g$:}
\example{rulesstate}
\fnentry{ContractionTable}{\fn{ContractionTable[state, \{i, j, \dots\}]} shows all the contractions between the modes of the
sites, $\langle c^\dagger_a c_b\rangle$ and $\langle c_ac_b\rangle$; with a third argument \texttt{"Normal"} or \texttt{"Anomalous"}
it gives the matrix, with the modes in the order of \texttt{Tuples[\{sites, \$Spins\}]}. The general state on one site, and
the symmetric one:}
\example{channels}
\example{symmetries}
\fnentry{Contraction}{\fn{Contraction[state, op1, op2]}; \fn{Contraction[op1, op2]} uses \fn{\$State}.}
\example{contraction}
\fnentry{SpinSingletState}{The usual shortcut: \texttt{GaussianState[\{"SpinRotation", "Charge"\}, "Filling" -> 1/2]}, or,
with \texttt{"Pairing" -> \textbackslash[CapitalDelta]}, the SU(2) state with singlet pairing and on-site pairing
\texttt{"OnSitePairing"} (default 0). Options \texttt{"Hopping"}, \texttt{"Filling"}, \texttt{"Pairing"},
\texttt{"OnSitePairing"}, \texttt{"Hermitian"}.}
\example{singletoptions}
\fnentry{GeneralState}{The general state written by spin components instead of channels: $\rho$ (option \texttt{"Normal"},
Hermitian) and, with \texttt{"Anomalous" -> \textbackslash[Kappa]}, $\kappa$, antisymmetric and written in canonical order.}
\example{generalstate}
\example{generalanomalous}
\fnentry{StateSelect, LocalPairQ, AnomalousPairQ}{\fn{StateSelect[state, crit]} keeps the contractions of the pairs with
\texttt{crit[op1, op2]} true. It is a filter, not a symmetry: it gives a physical state when the criterion is a symmetry or
a decoupling of the sites (\fn{LocalPairQ}: uncorrelated sites), and otherwise a way to isolate a class of terms.
\fn{AnomalousPairQ} is true for two $c$ or two $c^\dagger$.}
\example{stateselect}
\fnentry{\$State}{The state used when none is given, \texttt{SpinSingletState[]} by default.}

\section{Wick's theorem and mean field}
\fnentry{WickVEV}{\fn{WickVEV[expr, state]} is the expectation value by Wick's theorem: the expression is normal-ordered,
and each monomial is the sum over its pairings with the sign of the permutation. Options: \texttt{"MaxOrder" -> m}, at most
$m$ contractions between different sites; \texttt{"NormalOrder" -> False}, Wick on the expression as it is.}
\example{maxorder}
\fnentry{VerifyNormalOrder}{$0$ when \fn{WickVEV} with and without normal ordering agree.}
\example{verify}
\fnentry{MeanField}{\fn{MeanField[expr, state]} is the Hartree--Fock--Bogoliubov Hamiltonian of \texttt{expr}: every
normal-ordered product $O_1\cdots O_{2n}$ becomes
\[
\sum_{p<q}(-1)^{p+q+1}\,O_pO_q\,\Big\langle\prod_{r\ne p,q}O_r\Big\rangle-(n-1)\,\langle O_1\cdots O_{2n}\rangle ,
\]
two operators kept out in every way and the rest replaced by its expectation value. The result is quadratic, has the same
expectation value as \texttt{expr} and the same first variation with respect to the state: it is the Hamiltonian whose
Gibbs state makes $\langle H\rangle-TS$ stationary among Gaussian states. Quadratic terms are unchanged; odd terms give
zero. The density--density interaction (Hartree and Fock):}
\example{meanfield}
\example{meanfieldquadratic}

\section{\LaTeX}
\fnentry{ToTeX}{The \LaTeX{} of an expression: ladder and local operators, spin operators, amplitudes and spins.}
\example{totex}
\fnentry{WickPairings}{\fn{WickPairings[product, state]} lists the pairings of a product of ladder operators that do not vanish
in the state, in Wick's theorem applied to the product as written: \texttt{"Sign"}, \texttt{"Pairs"} (positions in the
product) and \texttt{"Value"} (the product of the contractions).}
\example{wickpairings}
\fnentry{WickTeX}{\fn{WickTeX[product, state]} writes Wick's theorem for a product in \LaTeX: the expectation value, the sum of
the pairings drawn as contractions, their values and the result. Options \texttt{"Values"} and \texttt{"Result"} (default
\texttt{True}) leave out the last two. The contractions use \pkg{simpler-wick}: \texttt{\textbackslash usepackage\{simpler-wick\}} in the
preamble, and two \LaTeX{} runs. Overlapping contractions are drawn at different heights.}
\example{wicktex}
\extypeset{wicktex}

\section{Gutzwiller projectors}
Loaded by \texttt{Needs["FermiWick`Gutzwiller`"]}.
\fnentry{GutzwillerProjector, \textbackslash[Theta]}{$P^G_i=\sqrt2\,[\cos\theta\,P^1_i+\sin\theta\,(P^0_i+D_i)]$, the half-filling
Gutzwiller operator with the variational angle $\theta$; $\theta=0$ projects on single occupation.}
\example{gutzwillerprojector}
\fnentry{LocalSandwich}{\fn{LocalSandwich[P, \{i, j, \dots\}, X]} is $P_iP_j\cdots X\cdots P_jP_i$, normal-ordered;
\fn{LocalSandwich[PL, sites, X, PR]} uses different operators on the two sides.}
\example{localsandwich}
\fnentry{GutzwillerVEV, GutzwillerSimplify}{\fn{GutzwillerVEV[expr, state]} is \fn{WickVEV} to first order in the contractions
between different sites (the hopping amplitudes of the state), simplified by \fn{GutzwillerSimplify}.}
\example{gutzwillervev}

\section{Lattices}
\fnentry{DefinePeierlsHopping, DefineFlux}{\fn{DefinePeierlsHopping[t, \textbackslash[CurlyPhi]]} makes
$t_{ab}=t\,e^{i\varphi_{ab}}$ with $\varphi_{ba}=-\varphi_{ab}$; \fn{DefineFlux[\textbackslash[CurlyPhi], \{i, j, k, \dots\}, \textbackslash[CapitalPhi]]}
fixes the sum of the phases around the oriented loop to $\Phi$, eliminating one free phase.}
\example{peierls}

\section{Matrices and exact diagonalisation}
Loaded by \texttt{Needs["FermiWick`Fock`"]}, in a cell of its own (the names must be known before the code that uses them
is read):
\example{fockload}
\textbf{Conventions.} The modes of a space are ladder operators of FermiWick, numbered site by site, then orbital by
orbital, the spin fastest: $\alpha=(\text{site}-1)\,n_o n_s+(\text{orbital}-1)\,n_s+\text{spin}$, the order of
\texttt{Tuples[\{sites, orbitals, spins\}]} and of \fn{ContractionTable}. A basis state is the string of the occupations
$(n_1,\dots,n_M)$, with index $k=\sum_\alpha n_\alpha 2^{M-\alpha}$ (the first mode is the most significant bit): the order of
\texttt{Tuples[\{0, 1\}, M]}. The Jordan--Wigner sign is $c_\alpha|n\rangle=(-1)^{\sum_{\beta<\alpha}n_\beta}|n-e_\alpha\rangle$. The
matrices are \texttt{SparseArray} objects; inside a sector the basis states are in increasing order of $k$.

\fnentry{FockSpace, FockSpaceQ}{\fn{FockSpace[sites]} is the space of spin-$\frac12$ fermions on the sites, with the modes
\fn{c[i, s]}; \fn{FockSpace[sites, "Orbitals" -> \{o1, \dots\}]} has the modes \fn{c[i, o, s]}, and \texttt{"Spins" -> None}
gives spinless modes \fn{c[i]} (or \fn{c[i, o]}); \texttt{"Spins" -> \{s1, s2\}} names the two spins, the first one up.
\texttt{space["Modes"]}, \texttt{space["Count"]} and \texttt{space["Dimension"]} ($2^M$) are its properties.}
\example{fockspace}
\fnentry{ModeIndex}{\fn{ModeIndex[space, c[l1, l2, \dots]]} (or \fn{cDag}) is the index of a mode;
\fn{ModeIndex[space, l1, l2, \dots]} the same from the labels. A mode outside the space gives a message and
\texttt{\$Failed}.}
\example{modeindex}
\fnentry{FockSector, FockSectorQ}{\fn{FockSector[space, "N" -> n, "Sz" -> sz, "Counts" -> \{\{modes, n\}, \dots\}]} keeps
the basis states with $n$ particles, the given $S^z$ (two spins) and $n$ particles in each set of modes (ladder operators
or indices). Every operator is built directly in a sector, without the whole space: sectors with fixed $N$ are enumerated
bit by bit and indexed through the combinatorial number system, so that time and memory scale with their dimension;
sectors without fixed $N$ need the whole space and are allowed up to \texttt{\$FockFullLimit} modes (24).
\fn{ClearFockCache[]} forgets the sectors built so far.}
\example{focksector}
\fnentry{FockMatrix}{\fn{FockMatrix[expr, target]} is the matrix of an expression of FermiWick: ladder operators, local and
composite operators, operators defined with \fn{DefineOperator} (also \fn{GutzwillerProjector}), products with
\texttt{**} in the order written (no normal order is needed), coefficients numbers or symbols. \texttt{target} is a space,
a sector (the image sector is found from the change of $N$, $S^z$ and of every count) or a pair
\texttt{\{fromSector, toSector\}}. The labels of every operator must be those of modes of the space; a symbolic spin
must be summed first (\texttt{Sum[\dots, \{s, \$Spins\}]}).}
\example{fockmatrix}
\example{fockmatrixsymbolic}
\fnentry{FockExpectation}{\fn{FockExpectation[expr, v, sector]} is $\langle v|\,\texttt{expr}\,|v\rangle$ for a normalized
vector \texttt{v} of the sector (or space); an operator that leaves the sector gives $0$. Here the singlet of two sites:}
\example{fockexpectation}
\fnentry{StateToKet}{\fn{StateToKet[target, v]} writes a vector as kets of the local occupations ($0$, $\uparrow$,
$\downarrow$, $\uparrow\downarrow$ for spin $\frac12$; the occupations of the modes otherwise), site by site.}
\example{statetoket}
\fnentry{OccupationBasis, BasisState}{\fn{OccupationBasis[target]} lists the occupation vectors of the basis states, in
order; \fn{BasisState[target, occupations]} is the basis vector with the given occupations.}
\example{basisstate}
\fnentry{SpinSquared}{\fn{SpinSquared[target]} is $\mathbf S^2$ of the whole space or of a sector (two spins), built from
$S^+$ and $S^z$: faster than the matrix of $\sum_{ij}\mathbf S_i\cdot\mathbf S_j$.}
\example{spinsquared}
\fnentry{LowestStates}{\fn{LowestStates[H, k]} gives \texttt{\{energies, vectors\}} of the $k$ lowest eigenstates of a
Hermitian matrix: dense up to dimension \texttt{\$FockDenseLimit} (3000), Arnoldi above, from several starting vectors
when $k>1$, so that degenerate multiplets are complete.}
\example{lowest}
\fnentry{LehmannGreen, GreenMatrix}{\fn{LehmannGreen[H, space, modes]} is the single-particle Green's function at zero
temperature, $G_{ab}(z)=\langle c_a(z-H+E_0)^{-1}c^\dagger_b\rangle+\langle c^\dagger_b(z+H-E_0)^{-1}c_a\rangle$, of a Hamiltonian that
conserves $N$ (a matrix on the whole space, or a function sector $\to$ matrix), for the modes given (ladder operators or
indices), averaged over the degenerate ground multiplet: an association with \texttt{"GroundEnergy"},
\texttt{"GroundSectors"}, \texttt{"Degeneracy"}, \texttt{"Poles"} and \texttt{"Residues"}. Options: \texttt{"N"} (the
particle numbers where the ground state is sought, \texttt{All}), \texttt{"UseSz"} (sectors of $S^z$ when $H$ conserves it,
\texttt{Automatic}), \texttt{"Tolerance"} ($10^{-9}$, the width of the ground multiplet), \texttt{"ResidueCutoff"}
($10^{-12}$). The residues add up to the identity (the last entry below). \fn{GreenMatrix[g, z]} evaluates it at a
complex frequency.}
\example{lehmann}
\example{greenmatrix}
\fnentry{ParticleHoleOperator, TimeReversal}{\fn{ParticleHoleOperator[space, first, last]} is the unitary $P$ with
$Pc_\alpha P^{-1}=c^\dagger_\alpha$ for the modes \texttt{first..last} (all by default): $\sigma^x$ on each of them times
$e^{-i\pi\sum_\alpha(\alpha-1)n_\alpha}$. \fn{TimeReversal[space]} gives \texttt{<|"U" -> U, "Apply" -> f, "Inverse" -> finv|>}
for the antiunitary $T=UK$ of spin-$\frac12$ fermions, with \texttt{f[v] = U.Conjugate[v]}.}
\example{particlehole}
\fnentry{SectorBlock}{\fn{SectorBlock[H, sector]} is the block of a matrix on the whole space that acts in the sector.}
\example{sectorblock}

\end{document}
